\documentclass[10pt, a4paper]{article}

\usepackage{latexsym}
\usepackage[empty]{fullpage}
\usepackage{titlesec}
\usepackage{marvosym}
\usepackage[usenames,dvipsnames]{color}
\usepackage{verbatim}
\usepackage{enumitem}
\usepackage{fancyhdr}
\usepackage[hidelinks]{hyperref}
\usepackage{xcolor}
\usepackage{xeCJK}
\usepackage{fontspec}

\setCJKmainfont{STSong}
\setCJKsansfont{STHeiti}

\pagestyle{fancy}
\fancyhf{}
\fancyfoot{}
\renewcommand{\headrulewidth}{0pt}
\renewcommand{\footrulewidth}{0pt}
\setlength{\footskip}{3.60004pt}

\addtolength{\oddsidemargin}{-0.5in}
\addtolength{\evensidemargin}{-0.5in}
\addtolength{\textwidth}{1.0in}
\addtolength{\topmargin}{-0.6in}
\addtolength{\textheight}{1.2in}

\urlstyle{same}

\raggedbottom
\raggedright
\setlength{\tabcolsep}{0in}

% Section formatting
\titleformat{\section}{
  \vspace{-1pt}\raggedright\Large\bfseries
}{}{0em}{}[\color{black}\titlerule \vspace{2pt}]

% Custom commands
\newcommand{\resumeItem}[1]{
  \item\small{
    {#1 \vspace{-1pt}}
  }
}
\newcommand{\resumeSubheading}[4]{
  \textbf{#1} \hfill \small{\textbf{#2}} \\
  \textit{\small#3} \hfill \small{\textit{#4}}
}

\newcommand{\resumeProjectHeading}[2]{
  \textbf{\small#1} \hfill \small{#2}
  \vspace{2pt}
}

\newcommand{\resumeSubheadingItalic}[4]{
  \textbf{\large #1} \hfill \small{#2} \\
  \small{#3} \hfill \small{#4}
  \vspace{2pt}
}

\newlist{resumeItemList}{itemize}{1}
\setlist[resumeItemList]{
  leftmargin=0.15in,
  label=$\bullet$,
  parsep=1pt,
  itemsep=2pt,
  topsep=2pt,
}

\begin{document}

%----------HEADER----------
\begin{center}
  {\LARGE \textbf{李易非}} \\ \vspace{4pt}
  \small \href{mailto:yifei28@illinois.edu}{yifei28@illinois.edu} $|$ +86 157-1120-2091 $|$ \href{https://github.com/yifei28}{github.com/yifei28} $|$ \href{https://www.linkedin.com/in/yifei-li-25535a326/}{LinkedIn: yifei-li-25535a326}
\end{center}
\vspace{2pt}

%-----------EDUCATION-----------
\section{教育经历}
\vspace{-2pt}
\begin{resumeItemList}
  \item \resumeSubheading
  {伊利诺伊大学厄巴纳-香槟分校 (UIUC)}{伊利诺伊州，香槟，美国}
  {计算机工程 理学学士}{2023.05 -- 2027.05}
  \vspace{2pt}
  \item \small{\textbf{GPA:} 3.95/4.0}
  \item \small{\textbf{核心课程：}数据结构、算法、数据库系统、分布式系统、并行编程、通信网络、计算机系统、LLM工程推理}
\end{resumeItemList}

%-----------SKILLS-----------
\section{技能}
\vspace{-2pt}
\begin{resumeItemList}
  \item \small{\textbf{编程语言：}Python, Java, JavaScript, C++, C, SQL}
  \item \small{\textbf{技术栈：}Docker, Node.js, Spring Boot, FastAPI, MySQL, Redis, CUDA, PyTorch, GCP, Nginx, Linux, Git, CI/CD}
  \item \small{\textbf{语言：}英语（流利）、普通话（母语）}
\end{resumeItemList}

%-----------EXPERIENCE-----------
\section{工作经历}
\vspace{-2pt}
\resumeSubheadingItalic
{海马云 $|$ 后端开发实习生}{北京，中国}
{2024.06 -- 2024.08}{}
\begin{resumeItemList}
  \resumeItem{分析CPU火焰图定位性能瓶颈，移除冗余日志调用，\textbf{CPU使用率降低16\%}。}
  \resumeItem{使用\textbf{Kibana}进行每周用户数据分析，监控系统性能指标，发现并推动优化方向。}
  \resumeItem{配合DevOps团队使用\textbf{Jenkins} CI/CD流水线完成灰度与全量发布，使用\textbf{Jira}跟踪任务。}
\end{resumeItemList}
\vspace{12pt}

\resumeSubheadingItalic
{科研助理 $|$ 张欢教授实验室}{伊利诺伊州，香槟，美国}
{2026.01 -- 至今}{}
\begin{resumeItemList}
  \resumeItem{参与\textbf{auto\_LiRPA}开源神经网络验证库：实现\textbf{BoundScatterElements}支持\texttt{scatter\_reduce}的可靠边界传播（区间分析 + CROWN），扩展\texttt{pow(x, y)}支持浮点指数。}
  \resumeItem{实现\textbf{PerturbationPolyhedra}，将库从$L_p$范数（box）扰动扩展至\textbf{任意多面体输入集合}（线性约束$Cx + d \leq 0$定义）；计算最小外包盒并传播约束收紧边界估计，支持\textbf{分支定界}子域分割验证。}
\end{resumeItemList}
\vspace{12pt}

\resumeSubheadingItalic
{课程助教 $|$ ECE408 并行计算}{伊利诺伊州，香槟，美国}
{2026.01 -- 至今}{}
\begin{resumeItemList}
  \resumeItem{改进CNN并行计算实验，添加\textbf{内核吞吐量指标}输出，将抽象描述改写为具体、可量化的优化目标。}
  \resumeItem{开发\textbf{特征图可视化工具}，调用学生GPU内核并生成HTML报告，展示各卷积层中间输出，辅助调试。}
  \resumeItem{每周答疑辅导学生调试并行计算作业，参与设计考试题目，监考判卷。}
\end{resumeItemList}

%-----------PROJECTS-----------
\section{项目经历}
\vspace{-2pt}

\resumeProjectHeading
{保安管理系统 $|$ \normalfont{\small{Spring Boot, Java, Python, FastAPI, MySQL, Redis, Docker}}}{2025.05 -- 2025.12}
\begin{resumeItemList}
  \resumeItem{构建后端系统，含\textbf{63个REST API}，跨17个控制器，支持统一分页、多参数过滤及双重认证（JWT + 微信OpenID）。}
  \resumeItem{设计13个JPA实体的关系模式，乐观锁并发控制，\textbf{25+自定义JPQL}配合JOIN FETCH消除N+1问题。}
  \resumeItem{实现\textbf{Redis ZSET延迟队列}（3队列），三层校验 + 分布式乐观锁，保证恰好一次执行。}
  \resumeItem{\textbf{端到端交付}：\textbf{Docker Compose}编排4个微服务，57个单元测试，GitHub Actions CI/CD自动化部署。}
  \resumeItem{独立开发\textbf{FastAPI人脸识别微服务}，\textbf{VGGFace2}嵌入（512维）+ 余弦相似度匹配，Spring Boot经multipart POST调用，用于签到验证。}
\end{resumeItemList}
\vspace{8pt}

\resumeProjectHeading
{RISC-V 操作系统内核 $|$ \normalfont{\small{C, RISC-V Assembly}}}{2025.03 -- 2025.05}
\begin{resumeItemList}
  \resumeItem{实现\textbf{Sv39三级页表}虚拟内存，支持进程地址空间隔离、写时复制和用户态缺页处理。}
  \resumeItem{构建抢占式调度器，支持\textbf{锁与条件变量}，实现fork/exec/exit/wait及\textbf{16个系统调用}。}
  \resumeItem{开发\textbf{VirtIO块设备驱动}（中断驱动I/O）及KTFS文件系统，支持直接/间接/双重间接块映射。}
  \resumeItem{实现\textbf{管道IPC}（环形缓冲区 + 读写同步），构建统一I/O抽象层支持设备、文件与内存流。}
\end{resumeItemList}
\vspace{8pt}

\resumeProjectHeading
{GPU优化CNN推理 $|$ \normalfont{\small{CUDA, C++}}}{2025.01 -- 2025.05}
\begin{resumeItemList}
  \resumeItem{实现\textbf{LeNet-5} GPU卷积前向内核，采用\textbf{共享内存分块}、im2col展开及融合卷积-GEMM计算。}
  \resumeItem{基于\textbf{Nsight Systems/Compute}性能分析，用\textbf{TF32 Tensor Core}加速GEMM，\textbf{CUDA流}重叠传输，推理\textbf{提速37\%}。}
  \resumeItem{评估\textbf{7种优化策略}（批量100--10000），通过吞吐量分析识别出内存带宽为大规模输入主要瓶颈。}
\end{resumeItemList}

\end{document}
